\scanpage{127}
\mspage{18}
that $S$ is the intersection of $R_{v_1},\ldots,R_{v_n}$ where $v_i$ is the real discrete valuation of $K$ with $R_{v_i}=S_{P_i}$, each $v_i$ has center $M$ in $R$ and $v_1,\ldots,v_n$ are the only valuations of $K$ whose valuation rings contain $R$. Furthermore, if $R$ is complete then $n=1$.

\uline{Proof}. Everything except the last statement follows from [A7, Proposition 3.18 of Section 13]. Now assume that $R$ is complete. Let $a$ be any element of $P_1$, let $f(X)=X^m+f_1X^{m-1}+\cdots+f_m$ be the monic polynomial of least degree with coefficients $f_i$ in $R$ such that $f(a)=0$, then $a\in P_1$ and $f_i\in R\subset S$ for $i=1,\ldots,m$ implies that $f_m\in P_1\cap R=P$; if some $f_i$ were not in $M$ then in view of the fact that $S$ has no non-zero zero divisors, by Hensel's Lemma [C\deleted{\unsure{0|O}}\inserted{o}, Theorem 4] applied to the monic polynomial $f(X)$ over the complete local domain $R$ we would get a monic polynomial $g(X)$ in $R[X]$ of degree smaller than $m$ for which $g(a)=0$, whence $f_i\in M$ for all $i$; now for any $j$, $P_j\supset P$ and hence $f_i\in P_j$ for all $i$ and consequently $a^m\in P_j$ and hence $a\in P_j$. Therefore $P_1\subset P_j$ for all $j$ and hence $n=1$.

LEMMA 5. Let $(R,M)$ be an $n$-dimensional regular local domain such that $R$ contains a coefficient field $k$ and let $s$ be a non negative integer. Then $R/M^{\unsure{S|s}}$ can in a natural way be considered to be a vector space over $k$ and then its $k$-dimension is $n(s)=\binom{n+s-1}{n}$, i.e., (i) there exist elements $y_1,\ldots,y_{n(s)}$ in $R$ such that if $c_1,\ldots,c_{n(s)}$ are elements in $k$ with
\mspage{126}
